\documentclass[10pt,a4paper]{amsart}
\usepackage[margin=19mm]{geometry}
\usepackage{amsmath,amssymb,mathtools}
\usepackage[scr=rsfs,cal=euler]{mathalfa}
\usepackage{xfrac}
\usepackage[unicode,hidelinks,bookmarks=false]{hyperref}
\newcommand{\ie}{i.e.\ }
\newcommand{\cf}{cf.\ }
\newcommand{\resp}{respectively\ }
\newcommand{\Subsection}{\subsection}
\providecommand{\nobreakdash}{\nobreak\hbox{-}}
\setlength{\emergencystretch}{2em}
\multlinegap0pt
\usepackage{algorithm}\usepackage{algpseudocode}
\usepackage{xcolor}
\usepackage{enumerate}
\usepackage[shortlabels]{enumitem}
\usepackage{amssymb}
\usepackage{cancel}
\usepackage{float}
\newtheorem{definition}{Definition}
\newtheorem{proposition}{Propositionn}
\newcommand{\T}{\mathcal}
\newcommand{\m}[1]{{\bf{#1}}}
\title{Coupled Tensor-Tensor Completion Method with Applications in Drug Repurposing}
\author{Maryam Bagherian, Albert Hung, Ivo Dinov, Joshua Welch}
\date{Source: arXiv:2609.03190v1 (CC BY 4.0)}
\begin{document}
\maketitle

\noindent\emph{Source and licence.} Coupled Tensor-Tensor Completion Method with Applications in Drug Repurposing. Version arXiv:2609.03190v1 (CC BY 4.0), distributed by arXiv under the Creative Commons Attribution 4.0 International licence, http://creativecommons.org/licenses/by/4.0/, as recorded in the arXiv licence metadata for this exact version and shown on the abstract page https://arxiv.org/abs/2609.03190v1. Full licence notice: \texttt{licenses/arxiv-2609.03190-CC-BY-4.0.txt}.\par\smallskip

\noindent\textit{Editorial scope.} Counted unit: the article's proposition prop:perturb (source0.tex lines 283--322). The article's complete original proof of it is delivered with it (source0.tex lines 323--393, one \texttt{proof} environment of 71 lines). No mathematics is written, repaired or re-derived: the statement and the proof are the article's own lines, in the article's own notation, and no retained line is substituted or re-flowed.  Retained uncounted as the source-local context the proof needs: lines 271--278.  Nothing was omitted from any retained span, so the package carries no omission record.  No retained line contains a citation, so the wrapper carries no bibliography and no bibliography entry.  No retained line contains a figure environment, an image-inclusion command or drawing code, so no image is delivered and the package declares no asset dependency.  Editorial formatting only, in addition: an AMS \texttt{amsart} preamble that restates the article's own declarations and macros used by the retained lines, namely the environments \texttt{definition}, \texttt{proposition} on separate counters, together with the standard packages those lines need.  The retained environments print in this document as Definition 1, Proposition 1 in order of appearance, whereas the article prints the same items with the numbering of its own full document.  The complete native source of the article as distributed in the arXiv e-print for this exact version is bundled unchanged as \texttt{sources/209274/source0.tex}.
\begin{definition}\label{def:maha}
Let $\T X\in \mathbb R^{n_1\times n_2\times \cdots\times n_K}$. Consider the multilinear transformation \\
$ \varphi:\mathbb R^{n_1\times n_2\times \cdots\times n_K}\to \mathbb R^{n_1\times n_2\times \cdots\times n_K}$, with
\begin{align}
\varphi(\T X)= \T X\times_1 \m L^{(1)}\times_2 \m L^{(2)}\cdots \times_K \m L^{(K)}, 
\end{align}
where the square matrices $\m L^{(\ell)}\in \mathbb R^{n_\ell\times n_\ell}$, for $\ell=1,\cdots, K$, are called the $\ell$-mode matrices. Here, if $\m L^{(\ell)}$, for $\ell=1, \cdots, K$ are orthogonal matrices, then $ \varphi$ recovers Euclidean distance.
\end{definition}

\begin{proposition}[Perturbation bound for the tensor Mahalanobis distance]
	\label{prop:perturb}
	Let $\T X,\T S,\T E\in
	\mathbb R^{n_1\times\cdots\times n_K}$, and let
	\[
	d_M(\T X,\T S)
	=
	\Big\|
	(\T X-\T S)
	\times_1\m L^{(1)}
	\cdots
	\times_K\m L^{(K)}
	\Big\|_F^2
	\]
	be the tensor Mahalanobis distance defined in
	Definition~\ref{def:maha}. Define
	\[
	C=\prod_{\ell=1}^K\|\m L^{(\ell)}\|_2.
	\]
	Then, for any perturbation $\T E$,
	\[
	\begin{aligned}
		&
		\left|
		d_M(\T X,\T S+\T E)
		-
		d_M(\T X,\T S)
		\right|
		\\
		&\le
		2C
		\sqrt{d_M(\T X,\T S)}
		\,\|\T E\|_F
		+
		C^2
		\|\T E\|_F^2.
	\end{aligned}
	\]
Consequently, on bounded subsets of the tensor space, the tensor Mahalanobis distance is locally Lipschitz continuous with respect to perturbations of the auxiliary tensor.
\end{proposition}

\begin{proof}
	Let
	\[
	\T A=
	(\T X-\T S)
	\times_1\m L^{(1)}
	\cdots
	\times_K\m L^{(K)},
	\]
	and
	\[
	\T B=
	\T E
	\times_1\m L^{(1)}
	\cdots
	\times_K\m L^{(K)}.
	\]
	Then
	\[
	d_M(\T X,\T S+\T E)
	=
	\|\T A-\T B\|_F^2,
	\]
	while
	\[
	\|\T A\|_F
	=
	\sqrt{d_M(\T X,\T S)}.
	\]
	
	Using the identity
	\[
	\bigl|
	\|\T A-\T B\|_F^2-\|\T A\|_F^2
	\bigr|
	\le
	2\|\T A\|_F\|\T B\|_F+\|\T B\|_F^2,
	\]
	together with the multilinear norm inequality
	\[
	\|\T B\|_F
	\le
	\prod_{\ell=1}^K
	\|\m L^{(\ell)}\|_2
	\,
	\|\T E\|_F
	=
	C\|\T E\|_F,
	\]
	we obtain
	\[
	\begin{aligned}
		\left|
		d_M(\T X,\T S+\T E)
		-
		d_M(\T X,\T S)
		\right|
		&=
		\bigl|
		\|\T A-\T B\|_F^2-\|\T A\|_F^2
		\bigr|\\
		&\le
		2\|\T A\|_F\|\T B\|_F+\|\T B\|_F^2\\
		&\le
		2C\sqrt{d_M(\T X,\T S)}\,\|\T E\|_F
		+
		C^2\|\T E\|_F^2,
	\end{aligned}
	\]
	which proves the desired perturbation estimate.
	\end{proof}

\end{document}
